\Needspace{12\baselineskip}
\section{Where the Memos Experiment Fits}
\label{sec:position}

\subsection{Recommended research claim}
The recommended framing is:
\begin{quote}
Under what conditions does deliberately prepared project context improve coding
agent development outcomes or reduce resource use beyond an existing documented
workflow, once preparation and context maintenance are included?
\end{quote}
This asks for evidence of benefit rather than a predetermined winner. A useful
result could be improved correctness, lower resources with acceptable
correctness, or lower maintenance effort with preserved context accuracy.
A negative or uncertain result would also answer the question within the tested
setting.

The main contribution would be a controlled within-project case study and
extension of existing evaluation approaches. It would not automatically be a
new general context architecture. The strongest practical distinction is the
comparison with \emph{existing project practice}, coupled with observed
preparation and upkeep. That proposed combination is supported by the evidence
review; it has not been established as unique across all literature.

\subsection{Why Memos supplies a meaningful baseline}
The pinned project documentation already includes agent instructions, a domain
terminology file, change-routing guidance and test commands
\citep{memos2026}. The existing-workflow condition therefore contains context.
Removing that guidance would change the question into a no-context ablation.
The intended comparison is whether additional preparation earns its place in a
project that already documents its workflow.

Memos also gives a concrete Go/TypeScript setting, rather than treating one
result on a particular benchmark as universal. Using another project and task
mix is a replication/extension opportunity, not proof that language differences
will produce an advantage. Downloaded source copies establish a common starting
revision; they do not establish a validated executable baseline.

\subsection{The most useful literature-to-experiment links}
\begin{enumerate}
\item \textbf{Test operational benefits.}
An evidence-based template is a preparation candidate \citep{shanto2026}.
Evaluate the resulting workflow on independent tasks; document coverage alone
is not acceptance evidence. Template adoption remains a proposal.
\item \textbf{Reconcile efficiency with correctness.}
Measure resources together with successful implementations, including failed
attempts. This allows a project-level test of whether apparently faster work
still satisfies requirements, rather than selecting only the outcome that
favors a method.
\item \textbf{Include context upkeep in the comparison.}
Observe shared code changes, resulting guidance defects, repairs and measured
update/review effort. This moves from descriptive edit histories toward a
bounded evaluation of maintenance work. A static task comparison alone cannot
answer the maintenance question.
\end{enumerate}

\subsection{How to justify a method without assuming it wins}
A defensible prepared-context method needs a proposed mechanism. For example,
accurate instructions could reduce repeated discovery of relevant commands or
constraints. The competing explanation is that those facts are already easy
to find in Memos, so the added material mainly consumes attention. This
mechanism can be tested through correctness, resource records and supporting
process observations.

As a preparation hypothesis, first audit native guidance for relevant omissions
or hard-to-find requirements. Record the source and intended purpose of every
addition. Avoid repeating already accessible information merely to fill a
heading; an absence of useful additions is itself informative about the
incremental question.

A large generic project summary is consequently a weak default intervention.
A stronger candidate is concise, project-specific guidance linked to verified
sources and task-relevant constraints. This is a hypothesis, not an established
best practice or a selected implementation. The comparison should document
what is added and how it is delivered, including whether the agent actually
loads it. Merely placing a file in a folder does not demonstrate exposure.

The third method should remain open. Its eventual design should name a failure
mechanism it intends to address and the evidence that would show improvement.
Neither retrieval nor a knowledge graph follows automatically from the review.
If it adds another mechanism, its preparation and updating costs must be part
of the same accounting boundary.

The preparation procedure itself can be an intervention, as the positive
refinement results of \citet{shepard2026} suggest. Their particular method is
not selected here. Any preparation or tuning must remain separate from scored
tasks. Versioned context and transfer are also studied by \citet{kerbel2026};
the contribution should therefore be the measured Memos comparison and upkeep
evidence, rather than a claim to have first proposed a project knowledge layer.

\subsection{What the experiment could support}
If a prepared workflow solves more of the prespecified tasks under the same
execution policy, with uncertainty narrow enough to support the contrast, the
study could report a correctness improvement for that task set and agent.
If correctness is shown to remain within an agreed acceptable range while
resources fall, it could support an efficiency claim. If context remains
accurate after source changes with less observed repair/review work, it could
support a maintenance-efficiency claim.

It could not establish universal superiority, organization-wide savings or a
causal effect of document structure alone. Complete workflows may deliberately
supply different information; their contrast combines content, presentation
and delivery. A later equal-information ablation would be needed to attribute
an effect specifically to representation or selection.
